\documentclass[12pt,a4paper]{article}

\usepackage[T2A]{fontenc}
\usepackage[utf8]{inputenc}
\usepackage[russian]{babel}
\usepackage{amsmath,amssymb,amsthm}
\usepackage{enumitem}
\usepackage{array}
\usepackage{longtable}
\usepackage{booktabs}
\usepackage{geometry}
\usepackage{xcolor}
\usepackage{listings}
\usepackage{hyperref}

\geometry{margin=2.2cm}

\hypersetup{
    colorlinks=true,
    linkcolor=blue,
    urlcolor=blue
}

\lstset{
    basicstyle=\ttfamily\small,
    breaklines=true,
    frame=single,
    columns=fullflexible,
    backgroundcolor=\color{gray!8},
    inputencoding=utf8,
    extendedchars=true,
    keepspaces=true,
    literate=
    {А}{{\CYRA}}1
    {Б}{{\CYRB}}1
    {В}{{\CYRV}}1
    {Г}{{\CYRG}}1
    {Д}{{\CYRD}}1
    {Е}{{\CYRE}}1
    {Ё}{{\CYRYO}}1
    {Ж}{{\CYRZH}}1
    {З}{{\CYRZ}}1
    {И}{{\CYRI}}1
    {Й}{{\CYRISHRT}}1
    {К}{{\CYRK}}1
    {Л}{{\CYRL}}1
    {М}{{\CYRM}}1
    {Н}{{\CYRN}}1
    {О}{{\CYRO}}1
    {П}{{\CYRP}}1
    {Р}{{\CYRR}}1
    {С}{{\CYRS}}1
    {Т}{{\CYRT}}1
    {У}{{\CYRU}}1
    {Ф}{{\CYRF}}1
    {Х}{{\CYRH}}1
    {Ц}{{\CYRC}}1
    {Ч}{{\CYRCH}}1
    {Ш}{{\CYRSH}}1
    {Щ}{{\CYRSHCH}}1
    {Ъ}{{\CYRHRDSN}}1
    {Ы}{{\CYRERY}}1
    {Ь}{{\CYRSFTSN}}1
    {Э}{{\CYREREV}}1
    {Ю}{{\CYRYU}}1
    {Я}{{\CYRYA}}1
    {а}{{\cyra}}1
    {б}{{\cyrb}}1
    {в}{{\cyrv}}1
    {г}{{\cyrg}}1
    {д}{{\cyrd}}1
    {е}{{\cyre}}1
    {ё}{{\cyryo}}1
    {ж}{{\cyrzh}}1
    {з}{{\cyrz}}1
    {и}{{\cyri}}1
    {й}{{\cyrishrt}}1
    {к}{{\cyrk}}1
    {л}{{\cyrl}}1
    {м}{{\cyrm}}1
    {н}{{\cyrn}}1
    {о}{{\cyro}}1
    {п}{{\cyrp}}1
    {р}{{\cyrr}}1
    {с}{{\cyrs}}1
    {т}{{\cyrt}}1
    {у}{{\cyru}}1
    {ф}{{\cyrf}}1
    {х}{{\cyrh}}1
    {ц}{{\cyrc}}1
    {ч}{{\cyrch}}1
    {ш}{{\cyrsh}}1
    {щ}{{\cyrshch}}1
    {ъ}{{\cyrhrdsn}}1
    {ы}{{\cyrery}}1
    {ь}{{\cyrsftsn}}1
    {э}{{\cyrerev}}1
    {ю}{{\cyryu}}1
    {я}{{\cyrya}}1
}

\newtheorem{definition}{Определение}
\newtheorem{principle}{Принцип}
\newtheorem{example}{Пример}
\newtheorem{remark}{Замечание}

\title{Билет №43 \\ \small Интерфейс пользователя (интерфейсы командной строки, текстовые интерфейсы, графические интерфейсы). Разработка пользовательского интерфейса. Мультимедийные среды интерфейсного взаимодействия. (ИТМО)}
\author{Сериков Владислав Александрович}
\date{13 мая 2026}

\begin{document}

\maketitle
\tableofcontents
\newpage

\section{Понятие пользовательского интерфейса}

\begin{definition}
\textbf{Пользовательский интерфейс} --- это совокупность средств, методов и правил взаимодействия пользователя с программной, аппаратной или программно-аппаратной системой.
\end{definition}

Пользовательский интерфейс включает:
\begin{itemize}
    \item способы ввода данных: клавиатура, мышь, сенсорный экран, голос, жесты, камера, датчики;
    \item способы вывода информации: текст, графика, звук, видео, вибрация, дополненная реальность;
    \item модель взаимодействия: команды, меню, формы, прямое манипулирование объектами, диалог, сценарии;
    \item правила обратной связи: сообщения об ошибках, индикаторы выполнения, уведомления;
    \item средства помощи пользователю: подсказки, документация, автодополнение, обучение внутри интерфейса.
\end{itemize}

В широком смысле пользовательский интерфейс является частью человеко-компьютерного взаимодействия
(\textit{Human--Computer Interaction, HCI}) --- области, изучающей проектирование, реализацию и оценку интерактивных вычислительных систем для использования человеком.

\section{Основные характеристики интерфейса}

Качество пользовательского интерфейса обычно оценивается по нескольким группам характеристик.

\subsection{Функциональные характеристики}

\begin{itemize}
    \item \textbf{Полнота}: интерфейс предоставляет доступ ко всем необходимым функциям системы.
    \item \textbf{Доступность функций}: пользователь может быстро найти нужное действие.
    \item \textbf{Соответствие предметной области}: термины, действия и структура интерфейса соответствуют реальным задачам пользователя.
\end{itemize}

\subsection{Эргономические характеристики}

\begin{itemize}
    \item \textbf{Удобство использования} --- насколько легко пользователь достигает своих целей.
    \item \textbf{Обучаемость} --- насколько быстро новый пользователь осваивает систему.
    \item \textbf{Запоминаемость} --- насколько легко пользователь возвращается к работе после перерыва.
    \item \textbf{Низкая частота ошибок} --- интерфейс предотвращает ошибки или помогает быстро их исправлять.
    \item \textbf{Удовлетворённость пользователя} --- субъективная оценка комфорта работы.
\end{itemize}

\subsection{Технические характеристики}

\begin{itemize}
    \item скорость отклика;
    \item надёжность;
    \item совместимость с устройствами;
    \item масштабируемость;
    \item локализация;
    \item поддержка доступности для пользователей с ограничениями зрения, слуха, моторики или когнитивных функций.
\end{itemize}

\section{Классификация пользовательских интерфейсов}

Пользовательские интерфейсы можно классифицировать по способу взаимодействия.

\begin{longtable}{p{4cm}p{5cm}p{5cm}}
\toprule
\textbf{Тип интерфейса} & \textbf{Особенности} & \textbf{Примеры} \\
\midrule
Командная строка & Взаимодействие через текстовые команды & Unix shell, PowerShell, SQL-консоль \\
Текстовый интерфейс & Экран состоит из текстовых окон, меню и форм & BIOS/UEFI setup, Midnight Commander, установщики ОС \\
Графический интерфейс & Используются окна, значки, меню, указатель & Windows, macOS, GNOME, мобильные приложения \\
Веб-интерфейс & Интерфейс доступен через браузер & Интернет-магазин, личный кабинет, облачные сервисы \\
Голосовой интерфейс & Команды задаются естественной речью & Голосовые помощники, IVR-системы \\
Жестовый интерфейс & Ввод осуществляется жестами тела или рук & Kinect, VR-контроллеры, сенсорные жесты \\
Мультимодальный интерфейс & Используются несколько каналов взаимодействия & Автомобильная информационная система, AR/VR-среда \\
\bottomrule
\end{longtable}

\section{Интерфейсы командной строки}

\subsection{Общее понятие}

\begin{definition}
\textbf{Интерфейс командной строки} --- это тип пользовательского интерфейса, в котором пользователь вводит текстовые команды, а система возвращает текстовый результат выполнения.
\end{definition}

Командная строка является одним из старейших и наиболее мощных видов интерфейса. Она широко используется администраторами, программистами, инженерами DevOps, специалистами по анализу данных и пользователями, которым требуется высокая степень контроля над системой.

\subsection{Структура команды}

Обычно команда имеет следующий вид:

\[
\texttt{command [options] [arguments]}
\]

где:
\begin{itemize}
    \item \texttt{command} --- имя программы или встроенной команды;
    \item \texttt{options} --- параметры, изменяющие поведение команды;
    \item \texttt{arguments} --- объекты, над которыми выполняется действие.
\end{itemize}

\begin{example}
Команда Unix:
\begin{lstlisting}
ls -la /home/user
\end{lstlisting}

Здесь:
\begin{itemize}
    \item \texttt{ls} --- команда вывода списка файлов;
    \item \texttt{-la} --- параметры: показать скрытые файлы и подробную информацию;
    \item \texttt{/home/user} --- каталог, содержимое которого требуется вывести.
\end{itemize}
\end{example}

\subsection{Преимущества командной строки}

\begin{itemize}
    \item высокая скорость работы для опытного пользователя;
    \item возможность автоматизации с помощью сценариев;
    \item низкие требования к вычислительным ресурсам;
    \item точный контроль над операциями;
    \item удобство удалённого администрирования;
    \item возможность композиции команд через конвейеры.
\end{itemize}

\subsection{Недостатки командной строки}

\begin{itemize}
    \item высокий порог входа;
    \item необходимость запоминать команды и параметры;
    \item низкая наглядность;
    \item риск разрушительных ошибок;
    \item слабая пригодность для задач, где важна визуальная манипуляция объектами.
\end{itemize}

\subsection{Конвейеры и композиция команд}

Одна из ключевых идей командной строки Unix --- композиция простых программ. Каждая программа выполняет одну задачу, а результат одной программы может передаваться другой программе.

\[
\texttt{program1 | program2 | program3}
\]

\begin{example}
Найти в файле строки со словом \texttt{error} и посчитать их количество:
\begin{lstlisting}
grep "error" app.log | wc -l
\end{lstlisting}

Шаги выполнения:
\begin{enumerate}
    \item \texttt{grep "error" app.log} выбирает строки, содержащие слово \texttt{error}.
    \item Символ \texttt{|} передаёт найденные строки следующей команде.
    \item \texttt{wc -l} подсчитывает количество строк.
\end{enumerate}
\end{example}

\subsection{Минимальный алгоритм обработки команды}

\textbf{Алгоритм интерпретации команды командной оболочкой:}

\begin{enumerate}
    \item Считать строку, введённую пользователем.
    \item Разбить строку на токены: имя команды, параметры, аргументы, операторы перенаправления.
    \item Выполнить подстановки: переменные окружения, шаблоны файлов, подстановку команд.
    \item Определить, является ли команда встроенной командой оболочки.
    \item Если команда встроенная, выполнить её внутри оболочки.
    \item Если команда внешняя, найти исполняемый файл в каталогах переменной \texttt{PATH}.
    \item Создать процесс для выполнения программы.
    \item Передать программе аргументы и окружение.
    \item Дождаться завершения или запустить процесс в фоне.
    \item Вернуть код завершения.
\end{enumerate}

\begin{example}
Пусть пользователь вводит:
\begin{lstlisting}
echo $HOME
\end{lstlisting}

Иллюстрация шагов:
\begin{enumerate}
    \item Считана строка \texttt{echo \$HOME}.
    \item Получены токены: \texttt{echo}, \texttt{\$HOME}.
    \item Переменная \texttt{\$HOME} заменяется, например, на \texttt{/home/student}.
    \item Обнаруживается встроенная или внешняя команда \texttt{echo}.
    \item Выполняется вывод:
\begin{lstlisting}
/home/student
\end{lstlisting}
\end{enumerate}
\end{example}

\section{Текстовые пользовательские интерфейсы}

\subsection{Общее понятие}

\begin{definition}
\textbf{Текстовый пользовательский интерфейс} --- это интерфейс, в котором экран представляется в виде текстовой сетки, а взаимодействие осуществляется через меню, формы, диалоговые окна и клавиатурные команды.
\end{definition}

Текстовый интерфейс занимает промежуточное положение между командной строкой и графическим интерфейсом. Пользователю уже не обязательно помнить точный синтаксис команд, но взаимодействие всё ещё происходит в текстовом режиме.

\subsection{Основные элементы текстового интерфейса}

\begin{itemize}
    \item строковые меню;
    \item псевдографические окна;
    \item таблицы;
    \item формы ввода;
    \item кнопки, изображённые текстом;
    \item горячие клавиши;
    \item строки состояния;
    \item диалоговые сообщения.
\end{itemize}

\begin{example}
Условная форма текстового интерфейса:
\begin{lstlisting}
+--------------------------------------+
|        Настройки пользователя         |
+--------------------------------------+
| Имя:       [ Ivan                 ]   |
| Почта:     [ ivan@example.com     ]   |
| Роль:      [ Пользователь       v ]   |
|                                      |
|        [ Сохранить ]  [ Отмена ]      |
+--------------------------------------+
\end{lstlisting}
\end{example}

\subsection{Преимущества текстовых интерфейсов}

\begin{itemize}
    \item работают без графической подсистемы;
    \item удобны для терминального доступа;
    \item менее требовательны к ресурсам;
    \item могут быть достаточно наглядными;
    \item позволяют использовать меню и формы;
    \item хорошо подходят для системного администрирования, установки ОС, встроенных систем.
\end{itemize}

\subsection{Недостатки текстовых интерфейсов}

\begin{itemize}
    \item ограниченные выразительные возможности;
    \item слабая поддержка сложной графики;
    \item зависимость от клавиатурной навигации;
    \item трудности с мультимедиа;
    \item меньшая привлекательность для массового пользователя по сравнению с GUI.
\end{itemize}

\section{Графические пользовательские интерфейсы}

\subsection{Общее понятие}

\begin{definition}
\textbf{Графический пользовательский интерфейс} --- это интерфейс, в котором пользователь взаимодействует с системой через визуальные объекты: окна, кнопки, меню, пиктограммы, панели, указатели, жесты и другие графические элементы.
\end{definition}

Графический интерфейс основан на идее визуального представления объектов и действий. Пользователь не столько вводит команды, сколько выбирает, перемещает и изменяет объекты на экране.

\subsection{WIMP-модель}

Классический графический интерфейс часто описывается моделью WIMP:

\[
\text{WIMP} = \text{Windows} + \text{Icons} + \text{Menus} + \text{Pointer}
\]

\begin{itemize}
    \item \textbf{Windows} --- окна, в которых отображаются документы и приложения.
    \item \textbf{Icons} --- значки объектов, файлов, действий.
    \item \textbf{Menus} --- меню команд.
    \item \textbf{Pointer} --- указатель мыши или аналогичный манипулятор.
\end{itemize}

\subsection{Прямое манипулирование}

\begin{definition}
\textbf{Прямое манипулирование} --- стиль взаимодействия, при котором пользователь видит объекты интерфейса и воздействует на них непосредственно, например перемещает файл мышью, изменяет размер окна или перетаскивает элемент.
\end{definition}

Основные свойства прямого манипулирования:
\begin{itemize}
    \item непрерывное представление объектов;
    \item быстрые, обратимые и инкрементальные действия;
    \item замена сложного синтаксиса команд физическими действиями;
    \item немедленная визуальная обратная связь.
\end{itemize}

\begin{example}
Удаление файла в графическом интерфейсе:
\begin{enumerate}
    \item Пользователь видит значок файла.
    \item Перетаскивает значок в корзину.
    \item Система подсвечивает корзину как допустимую цель.
    \item После отпускания кнопки мыши файл перемещается в корзину.
    \item Пользователь может отменить действие или восстановить файл.
\end{enumerate}
\end{example}

\subsection{Преимущества GUI}

\begin{itemize}
    \item высокая наглядность;
    \item низкий порог входа;
    \item поддержка визуального мышления;
    \item удобство работы с графикой, документами, мультимедиа;
    \item возможность непосредственного взаимодействия с объектами;
    \item хорошая поддержка обучения через исследование интерфейса.
\end{itemize}

\subsection{Недостатки GUI}

\begin{itemize}
    \item высокая ресурсоёмкость;
    \item сложность проектирования;
    \item трудность автоматизации некоторых действий;
    \item потенциальная перегруженность экрана;
    \item зависимость от качества визуального дизайна;
    \item для профессиональных пользователей GUI иногда медленнее командной строки.
\end{itemize}

\section{Сравнение CLI, TUI и GUI}

\begin{longtable}{p{3.3cm}p{3.7cm}p{3.7cm}p{3.7cm}}
\toprule
\textbf{Критерий} & \textbf{CLI} & \textbf{TUI} & \textbf{GUI} \\
\midrule
Основной ввод & Текстовые команды & Клавиатура, меню, формы & Мышь, клавиатура, сенсор, жесты \\
Наглядность & Низкая & Средняя & Высокая \\
Скорость для новичка & Низкая & Средняя & Высокая \\
Скорость для эксперта & Высокая & Средняя/высокая & Средняя \\
Автоматизация & Очень высокая & Средняя & Обычно ниже \\
Ресурсоёмкость & Низкая & Низкая/средняя & Средняя/высокая \\
Поддержка мультимедиа & Слабая & Слабая & Хорошая \\
Область применения & Администрирование, программирование & Установщики, серверные утилиты & Массовые приложения, офисные системы, веб \\
\bottomrule
\end{longtable}

\section{Разработка пользовательского интерфейса}

\subsection{Пользовательско-центрированный подход}

\begin{definition}
\textbf{Пользовательско-центрированное проектирование} --- это подход к разработке интерактивных систем, при котором потребности, ограничения, цели и контекст работы пользователей учитываются на всех этапах жизненного цикла системы.
\end{definition}

Основная идея: интерфейс проектируется не только исходя из внутренней структуры программы, но и исходя из задач пользователя.

\subsection{Основные этапы разработки интерфейса}

Разработка интерфейса обычно включает следующие этапы:

\begin{enumerate}
    \item Анализ пользователей и контекста использования.
    \item Выявление пользовательских задач.
    \item Формулирование требований к интерфейсу.
    \item Проектирование информационной архитектуры.
    \item Разработка сценариев взаимодействия.
    \item Создание прототипов.
    \item Визуальное проектирование.
    \item Реализация интерфейса.
    \item Тестирование удобства использования.
    \item Итерационное улучшение.
\end{enumerate}

\subsection{Алгоритм пользовательско-центрированного проектирования}

\textbf{Алгоритм:}

\begin{enumerate}
    \item Определить целевые группы пользователей.
    \item Описать цели пользователей.
    \item Описать контекст использования: устройства, среду, ограничения, частоту работы.
    \item Сформулировать пользовательские задачи.
    \item Определить требования к интерфейсу.
    \item Создать низкоуровневый прототип: бумажный, схематичный или интерактивный.
    \item Провести оценку прототипа с пользователями или экспертами.
    \item Выявить проблемы.
    \item Исправить проект.
    \item Повторять шаги 6--9 до достижения приемлемого качества.
    \item Реализовать интерфейс.
    \item Провести итоговое тестирование.
\end{enumerate}

\begin{example}
Разработка интерфейса приложения для записи к врачу:
\begin{enumerate}
    \item Пользователи: пациенты, администраторы клиники, врачи.
    \item Цель пациента: быстро выбрать специалиста и свободное время.
    \item Контекст: мобильный телефон, пользователь может находиться в дороге.
    \item Основная задача: ``записаться на приём''.
    \item Требования: крупные кнопки, минимум полей, понятный календарь.
    \item Прототип: экран выбора врача, экран выбора даты, экран подтверждения.
    \item Тестирование: пользователь должен записаться менее чем за 2 минуты.
    \item Исправления: добавить фильтр по ближайшему времени и кнопку ``повторить прошлую запись''.
\end{enumerate}
\end{example}

\section{Анализ пользователей и задач}

\subsection{Категории пользователей}

При проектировании интерфейса важно учитывать:
\begin{itemize}
    \item уровень компьютерной грамотности;
    \item профессиональные знания;
    \item возраст;
    \item физические возможности;
    \item мотивацию;
    \item частоту использования системы;
    \item уровень ответственности;
    \item условия работы.
\end{itemize}

Пользователей часто делят на:
\begin{itemize}
    \item новичков;
    \item периодических пользователей;
    \item опытных пользователей;
    \item экспертов.
\end{itemize}

\subsection{Персоны}

\begin{definition}
\textbf{Персона} --- это обобщённое описание типичного пользователя системы, включающее его цели, поведение, ограничения и контекст использования.
\end{definition}

\begin{example}
Персона:
\begin{quote}
Анна, 34 года, бухгалтер. Часто работает с отчётами, хорошо знает предметную область, но не любит сложные настройки. Цель --- быстро сформировать и отправить отчёт без ошибок.
\end{quote}
\end{example}

\subsection{Сценарии использования}

\begin{definition}
\textbf{Сценарий использования} --- это описание последовательности действий пользователя и системы для достижения конкретной цели.
\end{definition}

\begin{example}
Сценарий ``оплатить заказ'':
\begin{enumerate}
    \item Пользователь открывает корзину.
    \item Проверяет список товаров.
    \item Нажимает ``Оформить заказ''.
    \item Вводит адрес доставки.
    \item Выбирает способ оплаты.
    \item Подтверждает заказ.
    \item Получает сообщение об успешной оплате.
\end{enumerate}
\end{example}

\section{Принципы проектирования интерфейсов}

\subsection{Принцип видимости состояния системы}

\begin{principle}
Система должна своевременно информировать пользователя о своём состоянии с помощью понятной обратной связи.
\end{principle}

\begin{example}
Если файл загружается на сервер, интерфейс должен показывать индикатор прогресса, например:
\[
\text{Загрузка: } 65\%
\]
\end{example}

\subsection{Принцип соответствия реальному миру}

Интерфейс должен использовать понятные пользователю термины, а не внутренние технические обозначения системы.

\begin{example}
Лучше написать ``Корзина'', чем ``Временное хранилище удалённых объектов''.
\end{example}

\subsection{Принцип пользовательского контроля и свободы}

Пользователь должен иметь возможность отменять действия, возвращаться назад и исправлять ошибки.

\begin{example}
Кнопки:
\[
\text{Отменить}, \quad \text{Повторить}, \quad \text{Вернуться назад}
\]
снижают страх пользователя перед экспериментированием.
\end{example}

\subsection{Принцип согласованности}

Одинаковые действия и элементы должны выглядеть и работать одинаково во всей системе.

\begin{example}
Если синяя кнопка в одном окне означает ``Сохранить'', она не должна в другом окне означать ``Удалить''.
\end{example}

\subsection{Принцип предотвращения ошибок}

Лучше предотвратить ошибку, чем сообщать о ней после совершения.

\begin{example}
Если поле требует дату, можно использовать календарь вместо ручного ввода:
\[
\text{Дата: } [\text{выбрать из календаря}]
\]
\end{example}

\subsection{Принцип распознавания вместо вспоминания}

Интерфейс должен снижать нагрузку на память пользователя. Лучше показывать варианты, чем требовать запоминания команд.

\begin{example}
Меню ``Файл $\rightarrow$ Сохранить'' проще для новичка, чем необходимость помнить команду:
\begin{lstlisting}
save --current-document
\end{lstlisting}
\end{example}

\subsection{Принцип гибкости и эффективности}

Интерфейс должен быть удобен и для новичков, и для опытных пользователей.

\begin{example}
В текстовом редакторе:
\begin{itemize}
    \item новичок выбирает ``Файл $\rightarrow$ Сохранить'';
    \item опытный пользователь нажимает \texttt{Ctrl+S}.
\end{itemize}
\end{example}

\subsection{Принцип эстетичности и минимализма}

Интерфейс не должен содержать лишних элементов, отвлекающих от основной задачи.

\subsection{Принцип помощи в распознавании и исправлении ошибок}

Сообщения об ошибках должны быть понятными и конструктивными.

\begin{example}
Плохо:
\begin{quote}
Ошибка 0x80004005.
\end{quote}

Лучше:
\begin{quote}
Не удалось сохранить файл: нет прав записи в выбранную папку. Выберите другую папку или измените права доступа.
\end{quote}
\end{example}

\subsection{Принцип справки и документации}

Даже хороший интерфейс может требовать справки. Документация должна быть легко доступна, ориентирована на задачи пользователя и содержать примеры.

\section{Информационная архитектура интерфейса}

\begin{definition}
\textbf{Информационная архитектура} --- это структура организации, навигации и представления информации в интерфейсе.
\end{definition}

Она определяет:
\begin{itemize}
    \item какие разделы есть в системе;
    \item как они связаны;
    \item какие элементы находятся на каждом экране;
    \item как пользователь перемещается между разделами;
    \item какие данные показываются первыми.
\end{itemize}

\subsection{Основные схемы организации информации}

\begin{itemize}
    \item \textbf{Иерархическая}: разделы и подразделы.
    \item \textbf{Последовательная}: пользователь проходит шаги один за другим.
    \item \textbf{Сетевая}: переходы между множеством связанных объектов.
    \item \textbf{Поисковая}: пользователь находит объект через поиск и фильтры.
\end{itemize}

\begin{example}
Иерархия интернет-магазина:
\[
\text{Каталог} \rightarrow \text{Электроника} \rightarrow \text{Ноутбуки} \rightarrow \text{Игровые ноутбуки}
\]
\end{example}

\section{Прототипирование интерфейса}

\begin{definition}
\textbf{Прототип интерфейса} --- это предварительная модель интерфейса, используемая для проверки структуры, логики взаимодействия и визуальных решений до полной реализации.
\end{definition}

\subsection{Виды прототипов}

\begin{itemize}
    \item \textbf{Бумажный прототип}: экраны рисуются на бумаге.
    \item \textbf{Низкодетализированный прототип}: схематичные блоки без финального дизайна.
    \item \textbf{Высокодетализированный прототип}: близок к финальному интерфейсу.
    \item \textbf{Интерактивный прототип}: пользователь может нажимать элементы и переходить между экранами.
    \item \textbf{Вертикальный прототип}: реализована одна функция глубоко.
    \item \textbf{Горизонтальный прототип}: реализован внешний вид многих экранов без полной логики.
\end{itemize}

\subsection{Алгоритм бумажного тестирования прототипа}

\begin{enumerate}
    \item Подготовить бумажные экраны интерфейса.
    \item Сформулировать задачу для пользователя.
    \item Попросить пользователя выполнить задачу, ``нажимая'' на элементы.
    \item Модератор вручную заменяет бумажные экраны в ответ на действия пользователя.
    \item Наблюдатель фиксирует затруднения, ошибки и вопросы.
    \item После выполнения задачи проводится короткое интервью.
    \item Найденные проблемы группируются и приоритизируются.
    \item Прототип исправляется.
\end{enumerate}

\begin{example}
Задача: ``Найдите и оплатите курс по программированию''.

Если пользователь не замечает кнопку ``Купить'', значит:
\begin{itemize}
    \item кнопка расположена не там, где её ожидают;
    \item цвет или размер недостаточно выделяет действие;
    \item текст кнопки может быть непонятен.
\end{itemize}
\end{example}

\section{Оценка качества пользовательского интерфейса}

\subsection{Юзабилити-тестирование}

\begin{definition}
\textbf{Юзабилити-тестирование} --- это метод оценки интерфейса, при котором реальные или репрезентативные пользователи выполняют типовые задачи, а исследователь наблюдает за их действиями, ошибками и затруднениями.
\end{definition}

\subsection{Метрики юзабилити}

Основные метрики:
\begin{itemize}
    \item процент успешно выполненных задач;
    \item время выполнения задачи;
    \item количество ошибок;
    \item количество обращений к справке;
    \item субъективная удовлетворённость;
    \item число кликов или действий до достижения цели;
    \item коэффициент возвратов и отказов для веб-систем.
\end{itemize}

\subsection{Алгоритм юзабилити-тестирования}

\begin{enumerate}
    \item Определить цели тестирования.
    \item Выбрать целевые группы пользователей.
    \item Составить список задач.
    \item Подготовить прототип или рабочую версию системы.
    \item Провести инструктаж без подсказок по решению задач.
    \item Попросить пользователя выполнять задачи вслух, комментируя свои мысли.
    \item Зафиксировать ошибки, задержки, непонимание, обходные пути.
    \item Собрать субъективную оценку пользователя.
    \item Проанализировать результаты.
    \item Сформировать список проблем и рекомендаций.
\end{enumerate}

\begin{example}
Тестируется сайт библиотеки.

Задачи:
\begin{enumerate}
    \item Найти книгу по автору.
    \item Проверить, доступна ли она в электронном виде.
    \item Скачать библиографическое описание.
\end{enumerate}

Метрики:
\[
T_i = \text{время выполнения задачи } i,
\]
\[
S = \frac{\text{число успешно выполненных задач}}{\text{общее число задач}} \cdot 100\%.
\]
\end{example}

\subsection{Эвристическая оценка}

\begin{definition}
\textbf{Эвристическая оценка} --- это экспертный метод анализа интерфейса, при котором специалисты проверяют интерфейс на соответствие набору эвристик удобства использования.
\end{definition}

Преимущества:
\begin{itemize}
    \item быстро;
    \item относительно недорого;
    \item можно применять на ранних этапах;
    \item выявляет многие типовые проблемы.
\end{itemize}

Недостатки:
\begin{itemize}
    \item зависит от квалификации экспертов;
    \item не заменяет тестирование с реальными пользователями;
    \item может не выявить проблемы конкретного контекста использования.
\end{itemize}

\section{Диалоговые модели интерфейса}

\subsection{Понятие диалога}

\begin{definition}
\textbf{Диалог в пользовательском интерфейсе} --- это последовательность сообщений и действий между пользователем и системой, направленная на выполнение некоторой задачи.
\end{definition}

Диалог может быть:
\begin{itemize}
    \item командным;
    \item меню-ориентированным;
    \item формальным;
    \item графическим;
    \item голосовым;
    \item мультимодальным.
\end{itemize}

\subsection{Состояния интерфейса}

Многие интерфейсы удобно моделировать как конечный автомат:

\[
M = (S, A, \delta, s_0),
\]

где:
\begin{itemize}
    \item $S$ --- множество состояний интерфейса;
    \item $A$ --- множество действий пользователя;
    \item $\delta: S \times A \rightarrow S$ --- функция переходов;
    \item $s_0$ --- начальное состояние.
\end{itemize}

\begin{example}
Простой автомат формы входа:
\[
S = \{\text{пустая форма}, \text{данные введены}, \text{ошибка}, \text{успешный вход}\}.
\]

Действия:
\[
A = \{\text{ввести логин}, \text{ввести пароль}, \text{нажать вход}\}.
\]
\end{example}

\section{Обработка ошибок в интерфейсе}

Ошибки пользователя неизбежны, поэтому интерфейс должен:
\begin{itemize}
    \item предотвращать ошибки;
    \item обнаруживать ошибки;
    \item объяснять причины;
    \item предлагать исправление;
    \item сохранять данные пользователя;
    \item позволять отменять опасные действия.
\end{itemize}

\subsection{Типы ошибок}

\begin{itemize}
    \item \textbf{Ошибки slips}: пользователь знает цель, но случайно выполняет не то действие.
    \item \textbf{Ошибки mistakes}: пользователь неправильно понимает задачу или модель системы.
\end{itemize}

\begin{example}
\textbf{Slip}: пользователь случайно нажал ``Удалить'' вместо ``Переименовать''.

\textbf{Mistake}: пользователь думает, что кнопка ``Архивировать'' создаёт резервную копию, хотя она скрывает объект из основного списка.
\end{example}

\section{Доступность интерфейсов}

\begin{definition}
\textbf{Доступность интерфейса} --- это свойство системы быть пригодной для использования людьми с различными физическими, сенсорными, когнитивными и техническими ограничениями.
\end{definition}

Принципы доступности:
\begin{itemize}
    \item достаточный контраст текста и фона;
    \item возможность управления с клавиатуры;
    \item текстовые альтернативы изображениям;
    \item субтитры для аудио и видео;
    \item масштабируемый интерфейс;
    \item понятная структура заголовков;
    \item отсутствие зависимости только от цвета;
    \item совместимость со скринридерами.
\end{itemize}

\begin{example}
Плохо:
\[
\text{``Нажмите зелёную кнопку''}
\]

Лучше:
\[
\text{``Нажмите кнопку \textit{Подтвердить}''}
\]
поскольку пользователь с нарушением цветового восприятия может не различать цвет.
\end{example}

\section{Мультимедийные среды интерфейсного взаимодействия}

\subsection{Понятие мультимедийного интерфейса}

\begin{definition}
\textbf{Мультимедийный интерфейс} --- это интерфейс, в котором для взаимодействия используются различные формы представления информации: текст, графика, звук, видео, анимация, трёхмерная графика, тактильная обратная связь и другие каналы.
\end{definition}

Мультимедийные интерфейсы особенно важны в:
\begin{itemize}
    \item обучающих системах;
    \item компьютерных играх;
    \item медицинских тренажёрах;
    \item системах виртуальной и дополненной реальности;
    \item автомобильных системах;
    \item презентационных и справочных системах;
    \item системах дистанционного общения.
\end{itemize}

\subsection{Компоненты мультимедийного взаимодействия}

\begin{itemize}
    \item \textbf{Текст}: инструкции, подписи, сообщения.
    \item \textbf{Графика}: схемы, изображения, пиктограммы.
    \item \textbf{Анимация}: показ процессов во времени.
    \item \textbf{Аудио}: речь, сигналы, музыка.
    \item \textbf{Видео}: демонстрации, лекции, видеосвязь.
    \item \textbf{3D-графика}: пространственное моделирование.
    \item \textbf{Тактильная обратная связь}: вибрация, сопротивление, force feedback.
\end{itemize}

\subsection{Мультимодальность}

\begin{definition}
\textbf{Мультимодальный интерфейс} --- это интерфейс, который использует несколько модальностей ввода или вывода, например речь, жесты, взгляд, касание и визуальную обратную связь.
\end{definition}

\begin{example}
В автомобильной системе водитель может:
\begin{itemize}
    \item сказать голосом: ``Построить маршрут домой'';
    \item подтвердить выбор нажатием на экран;
    \item получить голосовую инструкцию;
    \item увидеть маршрут на карте.
\end{itemize}
\end{example}

\subsection{Преимущества мультимедийных интерфейсов}

\begin{itemize}
    \item высокая выразительность;
    \item возможность использовать разные каналы восприятия;
    \item повышение вовлечённости;
    \item удобство обучения через демонстрацию;
    \item поддержка пользователей с разными предпочтениями восприятия;
    \item реалистичность в симуляторах и тренажёрах.
\end{itemize}

\subsection{Недостатки мультимедийных интерфейсов}

\begin{itemize}
    \item высокая стоимость разработки;
    \item большие требования к ресурсам;
    \item сложность синхронизации медиа;
    \item риск когнитивной перегрузки;
    \item проблемы доступности;
    \item зависимость от оборудования и среды.
\end{itemize}

\section{Виртуальная и дополненная реальность}

\subsection{Виртуальная реальность}

\begin{definition}
\textbf{Виртуальная реальность} --- это компьютерно создаваемая интерактивная среда, в которой пользователь ощущает присутствие и может взаимодействовать с виртуальными объектами.
\end{definition}

Типичные устройства:
\begin{itemize}
    \item VR-шлем;
    \item контроллеры;
    \item системы отслеживания положения;
    \item перчатки с обратной связью;
    \item пространственный звук.
\end{itemize}

\subsection{Дополненная реальность}

\begin{definition}
\textbf{Дополненная реальность} --- это технология, при которой цифровые объекты накладываются на изображение реального мира и взаимодействуют с ним.
\end{definition}

\begin{example}
Приложение показывает через камеру смартфона, как мебель будет выглядеть в комнате.
\end{example}

\subsection{Особенности проектирования VR/AR-интерфейсов}

\begin{itemize}
    \item минимизация укачивания и сенсорного конфликта;
    \item естественные жесты;
    \item пространственное размещение элементов;
    \item учёт поля зрения;
    \item понятные точки фокуса;
    \item высокая частота кадров;
    \item безопасность физического пространства;
    \item корректная работа с глубиной и масштабом.
\end{itemize}

\section{Голосовые и разговорные интерфейсы}

\subsection{Голосовой интерфейс}

\begin{definition}
\textbf{Голосовой пользовательский интерфейс} --- это интерфейс, в котором основное взаимодействие пользователя с системой осуществляется с помощью речи.
\end{definition}

Голосовые интерфейсы применяются:
\begin{itemize}
    \item в умных колонках;
    \item телефонных справочных системах;
    \item автомобильных системах;
    \item мобильных помощниках;
    \item системах для людей с моторными ограничениями.
\end{itemize}

\subsection{Особенности голосового интерфейса}

Преимущества:
\begin{itemize}
    \item естественность для человека;
    \item возможность работы без рук;
    \item удобство при движении;
    \item помощь пользователям с ограничениями моторики или зрения.
\end{itemize}

Недостатки:
\begin{itemize}
    \item ошибки распознавания;
    \item шумовая зависимость;
    \item трудность исправления ошибок;
    \item слабая приватность;
    \item ограниченная видимость возможностей системы.
\end{itemize}

\subsection{Пример структуры голосового диалога}

\begin{example}
Пользователь:
\begin{quote}
Поставь будильник на семь утра.
\end{quote}

Система:
\begin{quote}
Будильник установлен на 7:00 завтра.
\end{quote}

Если команда неоднозначна:
\begin{quote}
На какую дату установить будильник?
\end{quote}
\end{example}

\section{Когнитивные аспекты пользовательского интерфейса}

\subsection{Ментальная модель}

\begin{definition}
\textbf{Ментальная модель} --- это внутреннее представление пользователя о том, как работает система.
\end{definition}

Хороший интерфейс согласует:
\begin{itemize}
    \item модель разработчика;
    \item модель системы;
    \item модель пользователя;
    \item визуальное представление интерфейса.
\end{itemize}

\subsection{Когнитивная нагрузка}

\begin{definition}
\textbf{Когнитивная нагрузка} --- это объём умственных усилий, необходимых пользователю для понимания интерфейса и выполнения задачи.
\end{definition}

Для снижения когнитивной нагрузки применяют:
\begin{itemize}
    \item группировку элементов;
    \item единообразие;
    \item визуальную иерархию;
    \item подсказки;
    \item автозаполнение;
    \item постепенное раскрытие сложности.
\end{itemize}

\section{Закон Фиттса как теоретическая основа интерфейсного проектирования}

В проектировании графических интерфейсов часто используется закон Фиттса, описывающий зависимость времени наведения на цель от расстояния до неё и размера цели.

\begin{definition}
\textbf{Закон Фиттса} в одной из распространённых форм записывается так:
\[
T = a + b \log_2\left(1 + \frac{D}{W}\right),
\]
где:
\begin{itemize}
    \item $T$ --- среднее время достижения цели;
    \item $D$ --- расстояние до цели;
    \item $W$ --- ширина цели в направлении движения;
    \item $a,b$ --- эмпирические коэффициенты.
\end{itemize}
\end{definition}

Из формулы следует:
\begin{itemize}
    \item чем дальше расположен элемент, тем дольше до него добираться;
    \item чем меньше элемент, тем сложнее попасть в него;
    \item крупные и часто используемые элементы следует размещать ближе к области внимания пользователя.
\end{itemize}

\begin{example}
Кнопка ``Отправить'' в форме должна быть достаточно крупной и находиться рядом с последним полем ввода. Маленькая кнопка в углу экрана увеличивает время наведения и вероятность ошибки.
\end{example}

\subsection{Практические следствия закона Фиттса}

\begin{itemize}
    \item важные кнопки должны быть достаточно большими;
    \item опасные действия не следует размещать слишком близко к часто используемым;
    \item края и углы экрана могут быть удобными целями, потому что указатель ограничен границами экрана;
    \item расстояние между связанными элементами должно быть небольшим.
\end{itemize}

\section{Закон Хика}

\begin{definition}
\textbf{Закон Хика} описывает зависимость времени выбора от количества альтернатив. В упрощённой форме:
\[
T = a + b \log_2(n + 1),
\]
где:
\begin{itemize}
    \item $T$ --- время принятия решения;
    \item $n$ --- число вариантов выбора;
    \item $a,b$ --- эмпирические коэффициенты.
\end{itemize}
\end{definition}

Практический смысл:
\begin{itemize}
    \item чем больше вариантов одновременно предъявлено пользователю, тем дольше он выбирает;
    \item длинные меню следует группировать;
    \item сложные формы лучше разбивать на шаги;
    \item редко используемые настройки можно скрывать в раздел ``Дополнительно''.
\end{itemize}

\begin{example}
Если пользователь выбирает категорию товара из 100 вариантов, лучше сгруппировать их:
\[
\text{Электроника} \rightarrow \text{Компьютеры} \rightarrow \text{Ноутбуки}
\]
чем показывать все 100 категорий одним списком.
\end{example}

\section{Визуальное проектирование интерфейса}

\subsection{Визуальная иерархия}

\begin{definition}
\textbf{Визуальная иерархия} --- это организация элементов интерфейса по степени важности с помощью размера, цвета, контраста, расположения и группировки.
\end{definition}

Средства визуальной иерархии:
\begin{itemize}
    \item размер;
    \item цвет;
    \item насыщенность;
    \item положение на экране;
    \item пустое пространство;
    \item шрифт;
    \item рамки и тени;
    \item анимация.
\end{itemize}

\subsection{Гештальт-принципы}

В интерфейсах часто используются психологические принципы группировки:

\begin{itemize}
    \item \textbf{близость}: элементы рядом воспринимаются как связанные;
    \item \textbf{сходство}: похожие элементы воспринимаются как группа;
    \item \textbf{замкнутость}: человек стремится видеть завершённые формы;
    \item \textbf{непрерывность}: взгляд следует по непрерывным линиям;
    \item \textbf{фигура и фон}: важный объект должен отделяться от фона.
\end{itemize}

\begin{example}
Поля ``Логин'' и ``Пароль'' следует расположить близко друг к другу, а кнопку ``Войти'' --- рядом с ними. Тогда пользователь воспринимает их как единую форму входа.
\end{example}

\section{Адаптивные и отзывчивые интерфейсы}

\begin{definition}
\textbf{Отзывчивый интерфейс} --- это интерфейс, который корректно изменяет расположение и размеры элементов в зависимости от размера экрана и устройства.
\end{definition}

\begin{definition}
\textbf{Адаптивный интерфейс} --- это интерфейс, который может изменять не только размеры, но и структуру, набор функций или сценарии взаимодействия в зависимости от пользователя, устройства или контекста.
\end{definition}

\begin{example}
На большом экране интернет-магазин показывает:
\begin{itemize}
    \item боковое меню категорий;
    \item сетку товаров в четыре колонки;
    \item фильтры слева.
\end{itemize}

На смартфоне:
\begin{itemize}
    \item меню скрывается в кнопку;
    \item товары идут одной колонкой;
    \item фильтры открываются в отдельной панели.
\end{itemize}
\end{example}

\section{Минимальный пример проектирования интерфейса}

Рассмотрим задачу: разработать интерфейс простого списка задач.

\subsection{Требования}

Пользователь должен уметь:
\begin{itemize}
    \item добавить задачу;
    \item отметить задачу выполненной;
    \item удалить задачу;
    \item просмотреть список задач.
\end{itemize}

\subsection{Вариант CLI}

\begin{lstlisting}
todo add "Подготовить конспект"
todo list
todo done 1
todo delete 1
\end{lstlisting}

Преимущества:
\begin{itemize}
    \item удобно автоматизировать;
    \item быстро для опытного пользователя.
\end{itemize}

Недостатки:
\begin{itemize}
    \item нужно помнить команды.
\end{itemize}

\subsection{Вариант TUI}

\begin{lstlisting}
+--------------------------+
|        Список задач       |
+--------------------------+
| [ ] Подготовить конспект  |
| [x] Купить книгу          |
|                          |
| [A] Добавить  [D] Удалить |
| [Space] Выполнено         |
+--------------------------+
\end{lstlisting}

Преимущества:
\begin{itemize}
    \item пользователь видит список;
    \item действия доступны через горячие клавиши.
\end{itemize}

\subsection{Вариант GUI}

\begin{lstlisting}
+--------------------------------+
| Мои задачи              [+]    |
+--------------------------------+
| [ ] Подготовить конспект       |
| [x] Купить книгу               |
| [ ] Повторить HCI              |
+--------------------------------+
\end{lstlisting}

Преимущества:
\begin{itemize}
    \item наглядность;
    \item работа мышью или касанием;
    \item легко добавить анимацию, фильтры, уведомления.
\end{itemize}

\section{Типовые ошибки при проектировании интерфейсов}

\begin{itemize}
    \item проектирование от структуры базы данных, а не от задач пользователя;
    \item перегрузка экрана элементами;
    \item отсутствие обратной связи;
    \item непонятные сообщения об ошибках;
    \item разные обозначения для одинаковых действий;
    \item слишком маленькие интерактивные элементы;
    \item отсутствие отмены опасных действий;
    \item скрытые основные функции;
    \item зависимость только от цвета;
    \item отсутствие тестирования с пользователями.
\end{itemize}

\section{Итоговые положения}

Пользовательский интерфейс является критически важной частью интерактивной системы. Командные интерфейсы обеспечивают мощность, точность и автоматизацию, но требуют высокой квалификации пользователя. Текстовые интерфейсы уменьшают необходимость запоминания команд и сохраняют низкие требования к ресурсам. Графические интерфейсы обеспечивают наглядность и удобство для широкого круга пользователей, но требуют более сложного проектирования.

Современная разработка интерфейсов строится на пользовательско-центрированном подходе, итерационном прототипировании, тестировании и учёте когнитивных возможностей человека. Мультимедийные и мультимодальные среды расширяют способы взаимодействия, используя звук, видео, жесты, речь, тактильную обратную связь, виртуальную и дополненную реальность. Однако их применение требует особого внимания к доступности, когнитивной нагрузке, техническим ограничениям и контексту использования.

\end{document}
